% =====================================================================
%  4. Justificación
% =====================================================================
\section{Justificación}
\label{sec:justificacion}

\subsection{Justificación académica: un espacio poco sintetizado y poco riguroso}
\label{sec:just-academica}

La sección~\ref{sec:contexto-aied} documenta dos hechos que, combinados, definen el
valor académico de esta investigación. El primero es de \emph{cobertura}: en la
revisión terciaria de \textcite{bond2024meta}, la evaluación de materiales de
aprendizaje con IA aparece en el 7,6\,\% de las revisiones y la curación de
materiales en el 4,5\,\%; en el mapeo previo de \textcite{zawacki2019systematic}, la
curación de materiales aparecía en 3 de 146 artículos. Cuando el tema aparece, se
aborda midiendo el uso de la plataforma, no el contenido.

El segundo es de \emph{rigor}: el 65\,\% de las revisiones analizadas por
\textcite{bond2024meta} tiene calidad metodológica de crítica a media, el 51,5\,\%
no reporta acuerdo entre revisores, y los autores piden explícitamente evaluar el
efecto de las herramientas sobre las personas y no solo su exactitud técnica. Diez
años antes, \textcite{zawacki2019systematic} ya señalaban estudios piloto con
resultados \enquote{preliminares positivos} que no controlan el efecto novedad.

La contribución académica que se pretende, entonces, no es \enquote{aplicar IA a la
educación}, sino producir sobre un caso concreto el tipo de evidencia que el campo
reclama: un benchmark con verdad de referencia anotada y acuerdo entre anotadores
reportado, configuraciones de comparación que permitan atribuir el desempeño a
componentes identificables, y un estudio con docentes que mida efecto sobre las
personas además de exactitud técnica.

\subsection{Justificación tecnológica: por qué esta combinación y no otra}
\label{sec:just-tecnologica}

Cada componente de la arquitectura responde a una limitación documentada, no a una
preferencia de implementación:

\begin{itemize}
  \item La \textbf{recuperación de evidencia} responde a que el criterio de
        corrección en este problema es la consistencia con el corpus del curso. Sin
        recuperación, el sistema dependería de la memoria paramétrica del modelo, que
        no es verificable contra los documentos del docente
        \parencite{lewis2020retrieval}.
  \item Los \textbf{detectores explícitos} responden a que la contradicción es la
        clase más difícil de las tareas de inferencia documental
        \parencite{koreeda2021contractnli,gubelmann2024capturing} y a que el
        aseguramiento automatizado disponible en repositorios educativos no llega al
        contenido, solo a metadatos y enlaces \parencite{clements2015open}.
  \item La \textbf{evidencia estructurada} responde a las causas de alucinación en
        sistemas RAG identificadas por \textcite{huang2025survey}, entre ellas las
        fuentes incorrectas o desactualizadas; la atribución explícita de origen es
        una de las mitigaciones que el propio survey señala.
  \item La \textbf{revisión humana} responde a que la evaluación automática por sí
        sola no basta para juzgar la sustancia de un recurso
        \parencite{tabares2017modelo,clements2015open}, y a que la revisión humana
        por sí sola subdetecta defectos \parencite{xie2018systematic}.
\end{itemize}

\begin{advertencia}
Que cada componente responda a una limitación documentada justifica \emph{su
inclusión en el diseño}; no demuestra que aporte valor en la práctica. Determinar
cuáles de estos componentes explican efectivamente el comportamiento del sistema es
precisamente el objeto de las ablaciones descritas en la Parte~V.
\end{advertencia}

\subsection{Justificación práctica: el riesgo de trasladar el costo al docente}
\label{sec:just-practica}

La justificación práctica de este proyecto no puede formularse como
\enquote{ahorrar tiempo al docente}, porque esa afirmación es exactamente la que
\textcite{su2026efficiency} muestran problemática. Se formula de otra manera: si un
asistente de curación propone cambios sin evidencia localizable, el docente debe
reconstruir por su cuenta la justificación de cada propuesta, y ese trabajo de
verificación es la forma específica de intensificación que la revisión sistemática
documenta.

El diseño que aquí se evalúa apuesta por reducir ese costo de verificación haciendo
que cada propuesta llegue acompañada del fragmento que la sustenta. La apuesta es
falsable: si el tiempo de revisión y la carga percibida no difieren de una condición
de revisión sin asistencia, la apuesta no se sostiene. El diseño experimental de la
Parte~V está construido para que ese resultado sea posible y reportable.

\subsection{Justificación normativa: qué exige un uso responsable}
\label{sec:just-normativa}

Dos marcos de referencia convergen en requisitos que este proyecto adopta como
restricciones de diseño y como objetos de medición.

El \emph{AI Risk Management Framework} del \textcite{nist2023airmf} sitúa la validez
y la fiabilidad como base de las demás características de una IA confiable, y
atraviesa todas ellas con transparencia y rendición de cuentas. Varias de sus
subcategorías se traducen directamente en obligaciones para este proyecto: documentar
los límites de conocimiento del sistema y cómo se supervisan sus salidas (MAP~2.2),
documentar los procesos de supervisión humana (MAP~3.5), documentar conjuntos de
prueba, métricas y herramientas de verificación (MEASURE~2.1), demostrar validez y
fiabilidad declarando los límites de generalización (MEASURE~2.5) y evaluar seguridad
y resiliencia (MEASURE~2.7). El marco recomienda además recoger datos sobre la
frecuencia y las razones con que los humanos anulan las salidas de la IA, lo que en
este proyecto se traduce en registrar los motivos de rechazo como dato de
investigación y no solo como funcionalidad.

La guía de \textcite{miao2023guidance} para IA generativa en educación e
investigación establece que la agencia humana no debe cederse en decisiones de alto
impacto, que las aplicaciones deben validarse ética y pedagógicamente antes de su
adopción institucional, y que la IA generativa, por rápida que sea como fuente de
información, no puede constituir una fuente autorizada de conocimiento
\parencite[p.~26]{miao2023guidance}. Sobre el apoyo al docente, la guía es explícita
en que los diseñadores humanos del currículo deben verificar el conocimiento factual
y la pertinencia de lo que la IA sugiera \parencite[p.~31]{miao2023guidance}.

\begin{estado}
\small
\textbf{Consecuencia de diseño.} Estos marcos no son un apartado ético decorativo:
determinan que la revisión humana previa a la publicación sea una restricción
arquitectónica no negociable, que la trazabilidad de cada decisión sea un
requerimiento y no una característica opcional, y que los límites de generalización
del sistema deban declararse junto con cualquier resultado que se reporte.
\end{estado}

\subsection{Justificación institucional}
\label{sec:just-institucional}

El proyecto se desarrolla en un contexto donde la evidencia sobre calidad de recursos
educativos tiene antecedentes locales: \textcite{tabares2017modelo} aplican su modelo
de evaluación por capas a la Federación de Repositorios de Objetos de Aprendizaje de
Colombia. Trabajar sobre corpus de cursos en español añade además una condición que
la literatura técnica revisada no cubre: la evidencia sobre recuperación y sobre
jueces automáticos está dominada por el inglés --- las tareas de recuperación de
MTEB son exclusivamente en inglés \parencite{muennighoff2023mteb} y los jueces
entrenados de ARES pierden concordancia con la anotación humana al cambiar de idioma
\parencite{saadfalcon2024ares} ---, lo que obliga a validar con datos en español los
componentes que se adopten. Esta restricción se desarrolla en la
Parte~V como una amenaza a la validez que el diseño debe controlar, no como un
supuesto que pueda darse por resuelto.
